% Replacement for lines 589-596 in Discussion section

This study builds on previous research into the selection and performance of centrality measures in LNA for identifying influential judgments. It addressed imbalanced data distributions and the phenomenon that certain centrality measures perform well at identifying relevant judgments (as defined by the ground truths) but not at finding less relevant judgments, whereas other measures do the opposite: they perform well identifying less relevant judgments but not relevant judgments.

In this study, Degree Centrality and PageRank emerged as particularly effective at identifying highly relevant judgments for both ground truths (Court Branch and Importance), whereas Degree Centrality, Eigenvector Centrality, and In-Degree Centrality demonstrated strong performance for identifying judgments of lower relevance. The findings led to the development of composite measures combining these complementary strengths, as a centrality measure ideally identifies both relevant and less relevant judgments.

Systematic evaluation of three composite combinations (PageRank + Degree, Degree + Eigenvector, Degree + In-Degree) revealed that all substantially outperformed individual centrality measures. The Degree + Eigenvector composite achieved the strongest performance overall, winning in up to 73.1\% of networks compared to 15-20\% for individual measures. This improvement demonstrates that leveraging distinct signals from high-relevance and low-relevance optimized centralities captures more comprehensive patterns of precedent value than any single measure.

The research further introduces balancing techniques, evaluates performance per ground-truth score (not just overall), and develops composite metrics through data-driven threshold optimization. Critically, the composite measures' consistent performance across balanced and unbalanced datasets indicates they capture generalizable structural characteristics of legal citation networks rather than artifacts of class imbalance.
